\part{Є҆ѵⷢ҇лїе ѿ Ма́рка}


% Chapter 1

\section % 1
\hskip0pt%
\St{З}ача́ло є҆ѵⷢ҇лїа і҆и҃са хрⷭ҇та̀, сн҃а бж҃їѧ,
\st{ꙗ҆́}коже є҆́сть пи́сано во прⷪ҇ро́цѣхъ: сѐ, а҆́зъ посыла́ю а҆́гг҃ла моего̀ пред̾ лице́мъ твои́мъ, и҆́же ᲂу҆гото́витъ пꙋ́ть тво́й пред̾ тобо́ю.
\st{Г}ла́съ вопїю́щагѡ въ пꙋсты́ни: ᲂу҆гото́вайте пꙋ́ть гдⷭ҇ень, пра̑вы твори́те стєзѝ є҆гѡ̀.
\st[]{Б}ы́сть і҆ѡа́ннъ крестѧ́й въ пꙋсты́ни и҆ проповѣ́даѧ кр҃ще́нїе покаѧ́нїѧ во ѿпꙋще́нїе грѣхѡ́въ.
\st{И҆} и҆схожда́ше къ немꙋ̀ всѧ̀ і҆ꙋде́йскаѧ страна̀ и҆ і҆ерⷭ҇ли́млѧне: и҆ креща́хꙋсѧ всѝ во і҆ѻрда́нѣ рѣцѣ̀ ѿ негѡ̀, и҆сповѣ́дающе грѣхѝ своѧ̑.
\st{Б}ѣ́ же і҆ѡа́ннъ ѡ҆болче́нъ власы̑ вельблꙋ̑жди, и҆ по́ѧсъ ᲂу҆сме́нъ ѡ҆ чре́слѣхъ є҆гѡ̀, и҆ ꙗ҆ды́й а҆крі̑ды и҆ ме́дъ ди́вїй.
\st{И҆} проповѣ́даше, глаго́лѧ: грѧде́тъ крѣ́плїй менє̀ в̾слѣ́дъ менє̀, є҆мꙋ́же нѣ́смь досто́инъ прекло́ньсѧ разрѣши́ти реме́нь сапѡ́гъ є҆гѡ̀:
\st{а҆́}зъ ᲂу҆́бѡ крести́хъ вы̀ водо́ю: то́й же крⷭ҇ти́тъ вы̀ дх҃омъ ст҃ы́мъ.
%
\section % 2
\st[]{И҆} бы́сть во ѻ҆́нѣхъ дне́хъ, прїи́де і҆и҃съ ѿ назаре́та галїле́йскагѡ и҆ крⷭ҇ти́сѧ ѿ і҆ѡа́нна во і҆ѻрда́нѣ.
\st{И҆} а҆́бїе восходѧ̀ ѿ воды̀, ви́дѣ разводѧ̑щасѧ небеса̀ и҆ дх҃а ꙗ҆́кѡ го́лꙋбѧ, сходѧ́ща на́нь.
\st{И҆} гла́съ бы́сть съ небесѐ: ты̀ є҆сѝ сн҃ъ мо́й возлю́бленный, ѡ҆ не́мже бл҃говоли́хъ.
\st{И҆} а҆́бїе дх҃ъ и҆зведѐ є҆го̀ въ пꙋсты́ню.
\st[]{И҆} бѣ̀ тꙋ̀ въ пꙋсты́ни дні́й четы́редесѧть, и҆скꙋша́емь сатано́ю, и҆ бѣ̀ со ѕвѣрьмѝ: и҆ а҆́гг҃ли слꙋжа́хꙋ є҆мꙋ̀.
\st{П}о преда́нїи же і҆ѡа́нновѣ, прїи́де і҆и҃съ въ галїле́ю, проповѣ́даѧ є҆ѵⷢ҇лїе црⷭ҇твїѧ бж҃їѧ
\st{и҆} гл҃ѧ, ꙗ҆́кѡ и҆спо́лнисѧ вре́мѧ и҆ прибли́жисѧ црⷭтвїе бж҃їе: пока́йтесѧ и҆ вѣ́рꙋйте во є҆ѵⷢ҇лїе.
%
\section % 3
\st[]{Х}одѧ́ же при мо́ри галїле́йстѣмъ, ви́дѣ сі́мѡна и҆ а҆ндре́а бра́та тогѡ̀ сі́мѡна, вмета̑юща мрє́жи въ мо́ре: бѣ́ста бо ры̑барѧ.
\st{И҆} речѐ и҆́ма і҆и҃съ: прїиди́та в̾слѣ́дъ менє̀, и҆ сотворю̀ ва́съ бы́ти ловца̑ человѣ́кѡмъ.
\st{И҆} а҆́бїе ѡ҆ста̑вльша мрє́жи своѧ̑, по не́мъ и҆до́ста.
\st{И҆} преше́дъ ма́лѡ ѿтꙋ́дꙋ, ᲂу҆зрѣ̀ і҆а́кѡва зеведе́ова и҆ і҆ѡа́нна бра́та є҆гѡ̀, и҆ та̑ въ кораблѝ стрѡ́ѧща мрє́жа:
\st{и҆} а҆́бїе воззва̀ ѧ҆̀. И҆ ѡ҆ста̑вльша ѻ҆тца̀ своего̀ зеведе́а въ кораблѝ съ нає́мники, по не́мъ и҆до́ста.
\st{И҆} внидо́ша въ капернаꙋ́мъ: и҆ а҆́бїе въ сꙋббѡ̑ты вше́дъ въ со́нмище, ᲂу҆ча́ше.
\st{И҆} дивлѧ́хꙋсѧ ѡ҆ ᲂу҆ч҃нїи є҆гѡ̀: бѣ́ бо ᲂу҆чѧ̀ и҆̀хъ ꙗ҆́кѡ вла́сть и҆мы́й, и҆ не ꙗ҆́кѡ кни́жницы.
%
\section % 4
\st[]{И҆} бѣ̀ въ со́нмищи и҆́хъ человѣ́къ въ дꙋ́сѣ нечи́стѣ, и҆ воззва̀,
\st{г}лаго́лѧ: ѡ҆ста́ви, что̀ на́мъ и҆ тебѣ̀, і҆и҃се назарѧни́не; прише́лъ є҆сѝ погꙋби́ти на́съ: вѣ́мъ тѧ̀, кто̀ є҆сѝ, ст҃ы́й бж҃їй.
\st{И҆} запретѝ є҆мꙋ̀ і҆и҃съ, гл҃ѧ: ᲂу҆молчѝ и҆ и҆зы́ди и҆з̾ негѡ̀.
\st{И҆} стрѧсѐ є҆го̀ дꙋ́хъ нечи́стый, и҆ возопѝ гла́сомъ вели́кимъ, и҆ и҆зы́де и҆з̾ негѡ̀.
\st{И҆} ᲂу҆жасо́шасѧ всѝ, ꙗ҆́коже стѧза́тисѧ и҆̀мъ къ себѣ̀, глаго́лющымъ: что̀ є҆́сть сїѐ; и҆ что̀ ᲂу҆ч҃нїе но́вое сїѐ, ꙗ҆́кѡ по ѡ҆́бласти и҆ дꙋховѡ́мъ нечи̑стымъ вели́тъ, и҆ послꙋ́шаютъ є҆гѡ̀;
\st{И҆}зы́де же слꙋ́хъ є҆гѡ̀ а҆́бїе во всю̀ странꙋ̀ галїле́йскꙋ.
%
\section % 5
\st[]{И҆} а҆́бїе и҆з̾ со́нмища и҆зше́дше, прїидо́ша въ до́мъ сі́мѡновъ и҆ а҆ндре́овъ со і҆а́кѡвомъ и҆ і҆ѡа́нномъ.
\st{Т}е́ща же сі́мѡнова лежа́ше ѻ҆гне́мъ жего́ма: и҆ а҆́бїе глаго́лаша є҆мꙋ̀ ѡ҆ не́й.
\st{И҆} пристꙋ́пль воздви́же ю҆̀, є҆́мь за рꙋ́кꙋ є҆ѧ̀: и҆ ѡ҆ста́ви ю҆̀ ѻ҆́гнь а҆́бїе, и҆ слꙋжа́ше и҆̀мъ.
\st{П}о́здѣ же бы́вшꙋ, є҆гда̀ захожда́ше со́лнце, приноша́хꙋ къ немꙋ̀ всѧ̑ недꙋ̑жныѧ и҆ бѣсны̑ѧ.
\st{И҆} бѣ̀ ве́сь гра́дъ собра́лсѧ къ две́ремъ.
\st{И҆} и҆сцѣлѝ мнѡ́ги ѕлѣ̀ стра́ждꙋщыѧ разли́чными недꙋ̑ги: и҆ бѣ́сы мнѡ́ги и҆згна̀, и҆ не ѡ҆ставлѧ́ше глаго́лати бѣ́сы, ꙗ҆́кѡ вѣ́дѧхꙋ є҆го̀ хрⷭ҇та̀ сꙋ́ща.
%
\linebreak
\section % 6
\st[]{И҆} ᲂу҆́трѡ, но́щи сꙋ́щей ѕѣлѡ̀, воста́въ и҆зы́де, и҆ и҆́де въ пꙋ́сто мѣ́сто, и҆ тꙋ̀ мл҃твꙋ дѣ́ѧше.
\st{И҆} гна́ша є҆го̀ сі́мѡнъ и҆ и҆̀же съ ни́мъ:
\st{и҆} ѡ҆брѣ́тше є҆го̀, глаго́лаша є҆мꙋ̀, ꙗ҆́кѡ всѝ тебѐ и҆́щꙋтъ.
\st{И҆} гл҃а и҆̀мъ: и҆́демъ въ бли̑жнїѧ вє́си и҆ гра́ды, да и҆ та́мѡ проповѣ́мъ: на сїе́ бо и҆зыдо́хъ.
\st{И҆} бѣ̀ проповѣ́даѧ на со́нмищихъ и҆́хъ, во все́й галїле́и, и҆ бѣ́сы и҆згонѧ̀.
\st[]{И҆} прїи́де къ немꙋ̀ прокаже́нъ, молѧ̀ є҆го̀ и҆ на колѣ̑нꙋ припа́даѧ пред̾ ни́мъ, и҆ глаго́лѧ є҆мꙋ̀, ꙗ҆́кѡ, а҆́ще хо́щеши, мо́жеши мѧ̀ ѡ҆чⷭ҇тити.
\st[]{І҆}и҃съ же млⷭ҇рдовавъ, просте́ръ рꙋ́кꙋ, коснꙋ́сѧ є҆гѡ̀, и҆ гл҃а є҆мꙋ̀: хощꙋ̀, ѡ҆чи́стисѧ.
\st{И҆} ре́кшꙋ є҆мꙋ̀, а҆́бїе ѿи́де ѿ негѡ̀ прокаже́нїе, и҆ чи́стъ бы́сть.
\st{И҆} запре́щь є҆мꙋ̀, а҆́бїе и҆згна̀ є҆го̀:
\st{и҆} гл҃а є҆мꙋ̀: блюдѝ, никомꙋ́же ничесѡ́же рцы̀: но ше́дъ покажи́сѧ і҆ере́еви и҆ принесѝ за ѡ҆чище́нїе твоѐ, ꙗ҆̀же повелѣ̀ мѡѷсе́й, во свидѣ́тельство и҆̀мъ.
\st{Ѻ҆́}нъ же и҆зше́дъ нача́тъ проповѣ́дати мно́гѡ и҆ проноси́ти сло́во, ꙗ҆́коже ктомꙋ̀ не мощѝ є҆мꙋ̀ ꙗ҆́вѣ во гра́дъ вни́ти: но внѣ̀ въ пꙋсты́хъ мѣ́стѣхъ бѣ̀. И҆ прихожда́хꙋ къ немꙋ̀ ѿвсю́дꙋ.

% Chapter 2

\vskip-4pt
\section % 7
\hskip0pt%
\St{И҆} вни́де па́ки въ капернаꙋ́мъ по дне́хъ: и҆ слы́шано бы́сть, ꙗ҆́кѡ въ домꙋ̀ є҆́сть.
\st{И҆} а҆́бїе собра́шасѧ мно́зи, ꙗ҆́коже ктомꙋ̀ не вмѣща́тисѧ ни при две́рехъ: и҆ гл҃аше и҆̀мъ сло́во.
\st{И҆} прїидо́ша къ немꙋ̀ носѧ́ще разсла́бленна жи́лами, носи́ма четы́рьми:
\st{и҆} не могꙋ́щымъ прибли́житисѧ къ немꙋ̀ наро́да ра́ди, ѿкры́ша покро́въ, и҆дѣ́же бѣ̀, и҆ прокопа́вше свѣ́сиша ѻ҆́дръ, на не́мже разсла́бленный лежа́ше.
\st[]{В}и́дѣвъ же і҆и҃съ вѣ́рꙋ и҆́хъ, гл҃а разсла́бленномꙋ: ча́до, ѿпꙋща́ютсѧ тебѣ̀ грѣсѝ твоѝ.
\st{Б}ѧ́хꙋ же нѣ́цыи ѿ кни̑жникъ тꙋ̀ сѣдѧ́ще и҆ помышлѧ́юще въ сердца́хъ свои́хъ:
\st{ч}то̀ се́й та́кѡ гл҃етъ хꙋлы̑; кто̀ мо́жетъ ѡ҆ставлѧ́ти грѣхѝ, то́кмѡ є҆ди́нъ бг҃ъ;
\st{И҆} а҆́бїе разꙋмѣ́въ і҆и҃съ дх҃омъ свои́мъ, ꙗ҆́кѡ та́кѡ ті́и помышлѧ́ютъ въ себѣ̀, речѐ и҆̀мъ: что̀ сїѧ̑ помышлѧ́ете въ сердца́хъ ва́шихъ;
\st{ч}то̀ є҆́сть ᲂу҆до́бѣе; рещѝ разсла́бленномꙋ: ѿпꙋща́ютсѧ тебѣ̀ грѣсѝ; и҆лѝ рещѝ: воста́ни, и҆ возмѝ ѻ҆́дръ тво́й, и҆ ходѝ;
\st[]{н}о да ᲂу҆вѣ́сте, ꙗ҆́кѡ вла́сть и҆́мать сн҃ъ чл҃вѣ́ческїй на землѝ ѿпꙋща́ти грѣхѝ: гл҃а разсла́бленномꙋ:
\st{т}ебѣ̀ гл҃ю: воста́ни, и҆ возмѝ ѻ҆́дръ тво́й, и҆ и҆дѝ въ до́мъ тво́й.
\st{И҆} воста̀ а҆́бїе, и҆ взе́мъ ѻ҆́дръ, и҆зы́де пред̾ всѣ́ми: ꙗ҆́кѡ диви́тисѧ всѣ̑мъ и҆ сла́вити бг҃а, глаго́лющымъ, ꙗ҆́кѡ николи́же та́кѡ ви́дѣхомъ.
%
\section % 8
\st[]{И҆} и҆зы́де па́ки къ мо́рю: и҆ ве́сь наро́дъ и҆дѧ́ше къ немꙋ̀, и҆ ᲂу҆ча́ше и҆̀хъ.
\st{И҆} мимогрѧды́й ви́дѣ леѵі́ю а҆лфе́ова, сѣдѧ́ща на мы́тницѣ, и҆ гл҃а є҆мꙋ̀: по мнѣ̀ грѧдѝ. И҆ воста́въ в̾слѣ́дъ є҆гѡ̀ и҆́де.
\st{И҆} бы́сть возлежа́щꙋ є҆мꙋ̀ въ домꙋ̀ є҆гѡ̀, и҆ мно́зи мытари̑ и҆ грѣ̑шницы возлежа́хꙋ со і҆и҃сомъ и҆ со ᲂу҆чн҃ки̑ є҆гѡ̀: бѧ́хꙋ бо мно́зи, и҆ по не́мъ и҆до́ша.
\st[]{И҆} кни́жницы и҆ фарїсе́є, ви́дѣвше є҆го̀ ꙗ҆дꙋ́ща съ мытари̑ и҆ грѣ̑шники, глаго́лахꙋ ᲂу҆чн҃кѡ́мъ є҆гѡ̀: что̀ ꙗ҆́кѡ съ мытари̑ и҆ грѣ̑шники ꙗ҆́стъ и҆ пїе́тъ;
\st{И҆} слы́шавъ і҆и҃съ гл҃а и҆̀мъ: не тре́бꙋютъ здра̑вїи врача̀, но болѧ́щїи: не прїидо́хъ призва́ти првⷣники, но грѣ́шники на покаѧ́нїе.
%
\section % 9
\st[]{И҆} бѧ́хꙋ ᲂу҆ченицы̀ і҆ѡа́ннѡвы и҆ фарїсе́йстїи постѧ́щесѧ. И҆ прїидо́ша и҆ глаго́лаша є҆мꙋ̀: почто̀ ᲂу҆ченицы̀ і҆ѡа́ннѡвы и҆ фарїсе́йстїи постѧ́тсѧ, а҆ твоѝ ᲂу҆чн҃цы̀ не постѧ́тсѧ;
\st{И҆} речѐ и҆̀мъ і҆и҃съ: є҆да̀ мо́гꙋтъ сы́нове бра́чнїи, до́ндеже жени́хъ съ ни́ми є҆́сть, пости́тисѧ; є҆ли́ко вре́мѧ съ собо́ю и҆́мꙋтъ жениха̀, не мо́гꙋтъ пости́тисѧ:
\st[]{п}рїи́дꙋтъ же дні́е, є҆гда̀ ѿи́метсѧ ѿ ни́хъ жени́хъ, и҆ тогда̀ постѧ́тсѧ въ ты̑ѧ дни̑:
\st{и҆} никто́же приложе́нїѧ пла́та небѣ́лена пришива́етъ къ ри́зѣ ве́тсѣ: а҆́ще ли же нѝ, во́зметъ коне́цъ є҆гѡ̀ но́вое ѿ ве́тхагѡ, и҆ го́рша дира̀ бꙋ́детъ:
\st{и҆} никто́же влива́етъ вїна̀ но́ва въ мѣ́хи вє́тхи: а҆́ще ли же нѝ, просади́тъ вїно̀ но́вое мѣ́хи, и҆ вїно̀ пролїе́тсѧ, и҆ мѣ́си поги́бнꙋтъ: но вїно̀ но́вое въ мѣ́хи нѡ́вы влїѧ́ти подоба́етъ.
%
\section % 10
\st[]{И҆} бы́сть мимоходи́ти є҆мꙋ̀ въ сꙋббѡ̑ты сквозѣ̀ сѣ̑ѧнїѧ, и҆ нача́ша ᲂу҆чн҃цы̀ є҆гѡ̀ пꙋ́ть твори́ти, востерза́юще кла́сы.
\st{И҆} фарїсе́є глаго́лахꙋ є҆мꙋ̀: ви́ждь, что̀ творѧ́тъ въ сꙋббѡ̑ты, є҆гѡ́же не досто́итъ;
\st{И҆} то́й гл҃аше и҆̀мъ: нѣ́сте ли николи́же члѝ, что̀ сотворѝ дв҃дъ, є҆гда̀ тре́бованїе и҆мѣ̀ и҆ взалка̀ са́мъ и҆ и҆̀же съ ни́мъ;
\st{к}а́кѡ вни́де въ до́мъ бж҃їй при а҆вїаѳа́рѣ а҆рхїере́и, и҆ хлѣ́бы предложе́нїѧ снѣдѐ, и҆́хже не досто́ѧше ꙗ҆́сти то́кмѡ і҆ере́ємъ, и҆ дадѐ и҆ сꙋ́щымъ съ ни́мъ;
\st[]{И҆} гл҃аше и҆̀мъ: сꙋббѡ́та человѣ́ка ра́ди бы́сть, а҆ не человѣ́къ сꙋббѡ́ты ра́ди:
\st{т}ѣ́мже госпо́дь є҆́сть сн҃ъ чл҃вѣ́ческїй и҆ сꙋббѡ́тѣ.

% Chapter 3

\vskip-4pt
\St{И҆} вни́де па́ки въ со́нмище: и҆ бѣ̀ та́мѡ человѣ́къ, сꙋ́хꙋ и҆мы́й рꙋ́кꙋ.
\st{И҆} назира́хꙋ є҆го̀, а҆́ще въ сꙋббѡ̑ты и҆сцѣли́тъ є҆го̀, да на́нь возглаго́лютъ.
\st{И҆} гл҃а человѣ́кꙋ сꙋ́хꙋ и҆мꙋ́щемꙋ рꙋ́кꙋ: ста́ни посредѣ̀.
\st{И҆} гл҃а и҆̀мъ: досто́итъ ли въ сꙋббѡ̑ты добро̀ твори́ти, и҆лѝ ѕло̀ твори́ти; дꙋ́шꙋ спастѝ, и҆лѝ погꙋби́ти; Ѻ҆ни́ же молча́хꙋ.
\st{И҆} воззрѣ́въ на ни́хъ со гнѣ́вомъ, скорбѧ̀ ѡ҆ ѡ҆камене́нїи серде́цъ и҆́хъ, гл҃а человѣ́кꙋ: прострѝ рꙋ́кꙋ твою̀. И҆ прострѐ: и҆ ᲂу҆тверди́сѧ рꙋка̀ є҆гѡ̀ цѣла̀ ꙗ҆́кѡ дрꙋга́ѧ.
%
\linebreak
\section % 11
\st[]{И҆} и҆зше́дше фарїсе́є, а҆́бїе со и҆рѡдїа̑ны совѣ́тъ творѧ́хꙋ на́нь, ка́кѡ є҆го̀ погꙋбѧ́тъ.
\st{І҆}и҃съ же ѿи́де со ᲂу҆чн҃ки̑ свои́ми къ мо́рю: и҆ мно́гъ наро́дъ ѿ галїле́и по не́мъ и҆́де, и҆ ѿ і҆ꙋде́и,
\st{и҆} ѿ і҆ерⷭ҇ли́ма, и҆ ѿ і҆дꙋме́и, и҆ со ѻ҆́нагѡ по́лꙋ і҆ѻрда́на. И҆ ѿ тѵ́ра и҆ сїдѡ́на мно́жество мно́гое, слы́шавше, є҆ли̑ка творѧ́ше, и҆ прїидо́ша къ немꙋ̀.
\st{И҆} речѐ ᲂу҆чн҃кѡ́мъ свои̑мъ, да кора́бль бꙋ́детъ ᲂу҆ негѡ̀ наро́да ра́ди, да не стꙋжа́ютъ є҆мꙋ̀.
\st[]{М}нѡ́ги бо и҆сцѣлѝ, ꙗ҆́коже напа́дати на него̀, да є҆мꙋ̀ прико́снꙋтсѧ, є҆ли́цы и҆мѣ́ѧхꙋ ра̑ны.
\st{И҆} дꙋ́си нечи́стїи, є҆гда̀ ви́дѧхꙋ є҆го̀, припа́дахꙋ къ немꙋ̀ и҆ зва́хꙋ, глаго́люще, ꙗ҆́кѡ ты̀ є҆сѝ сн҃ъ бж҃їй.
\st{И҆} мно́гѡ преща́ше и҆̀мъ, да не ꙗ҆вле́на є҆го̀ сотворѧ́тъ.
%
\section % 12
\st[]{И҆} взы́де на горꙋ̀ и҆ призва̀, и҆̀хже хотѧ́ше са́мъ: и҆ прїидо́ша къ немꙋ̀.
\st{И҆} сотворѝ двана́десѧте, да бꙋ́дꙋтъ съ ни́мъ, и҆ да посыла́етъ и҆̀хъ проповѣ́дати,
\st{и҆} и҆мѣ́ти вла́сть цѣли́ти недꙋ́ги и҆ и҆згони́ти бѣ́сы:
\st[]{и҆} наречѐ сі́мѡнꙋ и҆́мѧ пе́тръ:
\st{и҆} і҆а́кѡва зеведе́ова и҆ і҆ѡа́нна бра́та і҆а́кѡвлѧ: и҆ наречѐ и҆́ма и҆мена̀ воанерге́съ, є҆́же є҆́сть сы̑на гро́мѡва:
\st{и҆} а҆ндре́а, и҆ фїлі́ппа, и҆ варѳоломе́а, и҆ матѳе́а, и҆ ѳѡмꙋ̀, и҆ і҆а́кѡва а҆лфе́ова, и҆ ѳадде́а, и҆ сі́мѡна канані́та,
\st{и҆} і҆ꙋ́дꙋ і҆скарїѡ́тскаго, и҆́же и҆ предадѐ є҆го̀.
%
\section % 13
\st[]{И҆} прїидо́ша въ до́мъ: и҆ собра́сѧ па́ки наро́дъ, ꙗ҆́кѡ не мощѝ и҆̀мъ ни хлѣ́ба ꙗ҆́сти.
\st{И҆} слы́шавше и҆̀же бѧ́хꙋ ᲂу҆ негѡ̀, и҆зыдо́ша, да и҆́мꙋтъ є҆го̀: глаго́лахꙋ бо, ꙗ҆́кѡ неи́стовъ є҆́сть.
\st{И҆} кни́жницы, и҆̀же ѿ і҆ерⷭ҇ли́ма низше́дшїи, глаго́лахꙋ, ꙗ҆́кѡ веельзевꙋ́ла и҆́мать и҆ ꙗ҆́кѡ ѡ҆ кнѧ́зи бѣсо́встѣмъ и҆зго́нитъ бѣ́сы.
\st[]{И҆} призва́въ и҆̀хъ, въ при́тчахъ гл҃аше и҆̀мъ: ка́кѡ мо́жетъ сатана̀ сатанꙋ̀ и҆згони́ти;
\st{И҆} а҆́ще ца́рство на сѧ̀ раздѣли́тсѧ, не мо́жетъ ста́ти ца́рство то̀:
\st{и҆} а҆́ще до́мъ на сѧ̀ раздѣли́тсѧ, не мо́жетъ ста́ти до́мъ то́й:
\st{и҆} а҆́ще сатана̀ воста̀ на сѧ̀ са́мъ и҆ раздѣли́сѧ, не мо́жетъ ста́ти, но коне́цъ и҆́мать.
\st{Н}икто́же мо́жетъ сосꙋ́ды крѣ́пкагѡ, вше́дъ въ до́мъ є҆гѡ̀, расхи́тити, а҆́ще не пе́рвѣе крѣ́пкаго свѧ́жетъ: и҆ тогда̀ до́мъ є҆гѡ̀ расхи́титъ.
%
\section % 14
\st[]{А҆}ми́нь гл҃ю ва́мъ, ꙗ҆́кѡ всѧ̑ ѿпꙋ́стѧтсѧ согрѣшє́нїѧ сынѡ́мъ человѣ́чєскимъ, и҆ хꙋлє́нїѧ, є҆ли̑ка а҆́ще восхꙋ́лѧтъ:
\st{а҆} и҆́же восхꙋ́литъ на дх҃а ст҃а́го, не и҆́мать ѿпꙋще́нїѧ во вѣ́ки, но пови́ненъ є҆́сть вѣ́чномꙋ сꙋдꙋ̀.
\st{З}анѐ глаго́лахꙋ: дꙋ́ха нечи́стаго и҆́мать.
\st[]{П}рїидо́ша ᲂу҆́бѡ мт҃и и҆ бра́тїѧ є҆гѡ̀, и҆ внѣ̀ стоѧ́ще посла́ша къ немꙋ̀, зовꙋ́ще є҆го̀.
\st{И҆} сѣдѧ́ше наро́дъ ѡ҆́крестъ є҆гѡ̀. Рѣ́ша же є҆мꙋ̀: сѐ, мт҃и твоѧ̀ и҆ бра́тїѧ твоѧ̑ и҆ сєстры̀ твоѧ̑ внѣ̀ и҆́щꙋтъ тебѐ.
\st{И҆} ѿвѣща̀ и҆̀мъ, гл҃ѧ: кто̀ є҆́сть мт҃и моѧ̀, и҆лѝ бра́тїѧ моѧ̑;
\st[]{И҆} соглѧ́давъ ѡ҆́крестъ себє̀ сѣдѧ́щыѧ, гл҃а: сѐ, мт҃и моѧ̀ и҆ бра́тїѧ моѧ̑:
\st{и҆́}же бо а҆́ще сотвори́тъ во́лю бж҃їю, се́й бра́тъ мо́й и҆ сестра̀ моѧ̀ и҆ мт҃и мѝ є҆́сть.

% Chapter 4

\vskip-4pt
\section % 15
\hskip0pt%
\St{И҆} па́ки нача́тъ ᲂу҆чи́ти при мо́ри: и҆ собра́сѧ къ немꙋ̀ наро́дъ мно́гъ, ꙗ҆́коже самомꙋ̀ влѣ́зшꙋ въ кора́бль, сѣдѣ́ти въ мо́ри: и҆ ве́сь наро́дъ при мо́ри на землѝ бѧ́ше.
\st{И҆} ᲂу҆ча́ше и҆̀хъ при́тчами мно́гѡ и҆ гл҃аше и҆̀мъ во ᲂу҆ч҃нїи свое́мъ:
\st{с}лы́шите: сѐ, и҆зы́де сѣ́ѧй сѣ́ѧти:
\st[]{и҆} бы́сть є҆гда̀ сѣ́ѧше, ѻ҆́во падѐ при пꙋтѝ, и҆ прїидо́ша пти̑цы, и҆ позоба́ша є҆̀:
\st{д}рꙋго́е же падѐ при ка́мени, и҆дѣ́же не и҆мѧ́ше землѝ мно́ги, и҆ а҆́бїе прозѧбѐ, занѐ не и҆мѧ́ше глꙋбины̀ земны́ѧ:
\st{с}о́лнцꙋ же возсїѧ́вшꙋ присвѧ́де, и҆ занѐ не и҆мѧ́ше ко́рене, и҆́зсше:
\st{и҆} дрꙋго́е падѐ въ те́рнїи, и҆ взы́де те́рнїе, и҆ подавѝ є҆̀: и҆ плода̀ не дадѐ:
\st{и҆} дрꙋго́е падѐ на землѝ до́брѣй, и҆ даѧ́ше пло́дъ восходѧ́щь и҆ растꙋ́щь, и҆ припло́доваше на три́десѧть, и҆ на шестьдесѧ́тъ, и҆ на сто̀.
\st{И҆} гл҃аше: и҆мѣ́ѧй ᲂу҆́шы слы́шати да слы́шитъ.
%
\linebreak
\section % 16
\st[]{Є҆}гда́ же бы́сть є҆ди́нъ, вопроси́ша є҆го̀, и҆̀же бѧ́хꙋ съ ни́мъ, со ѻ҆бѣмана́десѧте ѡ҆ при́тчи.
\st{И҆} гл҃аше и҆̀мъ: ва́мъ є҆́сть дано̀ вѣ́дати та̑йны црⷭ҇твїѧ бж҃їѧ: ѡ҆́нѣмъ же внѣ̑шнимъ въ при́тчахъ всѧ̑ быва́ютъ,
\st{д}а ви́дѧще ви́дѧтъ, и҆ не ᲂу҆́зрѧтъ: и҆ слы́шаще слы́шатъ, и҆ не разꙋмѣ́ютъ: да не когда̀ ѡ҆братѧ́тсѧ, и҆ ѡ҆ста́вѧтсѧ и҆̀мъ грѣсѝ.
\st{И҆} гл҃а и҆̀мъ: не вѣ́сте ли при́тчи сеѧ̀; и҆ ка́кѡ всѧ̑ при̑тчи ᲂу҆разꙋмѣ́ете;
\st[]{С}ѣ́ѧй, сло́во сѣ́етъ.
\st{С}і́и же сꙋ́ть, и҆̀же при пꙋтѝ, и҆дѣ́же сѣ́етсѧ сло́во, и҆ є҆гда̀ ᲂу҆слы́шатъ, а҆́бїе прихо́дитъ сатана̀ и҆ ѿе́млетъ сло́во сѣ́ѧнное въ сердца́хъ и҆́хъ.
\st{И҆} сі́и сꙋ́ть та́кожде и҆̀же на ка́менныхъ сѣ́емїи, и҆̀же є҆гда̀ ᲂу҆слы́шатъ сло́во, а҆́бїе съ ра́достїю прїе́млютъ є҆̀:
\st{и҆} не и҆́мꙋтъ коре́нїѧ въ себѣ̀, но привре́менни сꙋ́ть: та́же бы́вшей печа́ли и҆лѝ гоне́нїю словесѐ ра́ди, а҆́бїе соблажнѧ́ютсѧ.
\st{А҆} сі́и сꙋ́ть, и҆̀же въ те́рнїи сѣ́емїи, слы́шащїи сло́во:
\st{и҆} печа̑ли вѣ́ка сегѡ̀, и҆ ле́сть бога́тства, и҆ ѡ҆ про́чихъ по́хѡти входѧ́щыѧ подавлѧ́ютъ сло́во, и҆ безпло́дно быва́етъ.
\st{И҆} сі́и сꙋ́ть, и҆̀же на землѝ до́брѣй сѣ́ѧннїи, и҆̀же слы́шатъ сло́во и҆ прїе́млютъ, и҆ пло́дствꙋютъ на три́десѧть, и҆ на шестьдесѧ́тъ, и҆ на сто̀.
\st[]{И҆} гл҃аше и҆̀мъ: є҆да̀ свѣти́льникъ прихо́дитъ, да под̾ спꙋ́домъ положа́тъ є҆го̀ и҆лѝ под̾ ѻ҆дро́мъ; не да ли на свѣ́щницѣ положе́нъ бꙋ́детъ;
\st{н}ѣ́сть бо та́йно, є҆́же не ꙗ҆ви́тсѧ, нижѐ бы́сть потае́но, но да прїи́детъ въ ꙗ҆вле́нїе:
\st{а҆́}ще кто̀ и҆́мать ᲂу҆́шы слы́шати, да слы́шитъ.
%
\linebreak
\section % 17
\st[]{И҆} гл҃аше и҆̀мъ: блюди́те что̀ слы́шите: въ ню́же мѣ́рꙋ мѣ́рите, возмѣ́ритсѧ ва́мъ, и҆ приложи́тсѧ ва́мъ слы́шащымъ:
\st{и҆́}же бо а҆́ще и҆́мать, да́стсѧ є҆мꙋ̀: а҆ и҆́же не и҆́мать, и҆ є҆́же и҆́мать, ѿи́метсѧ ѿ негѡ̀.
\st{И҆} гл҃аше: та́кѡ є҆́сть и҆ црⷭ҇твїе бж҃їе, ꙗ҆́коже человѣ́къ вмета́етъ сѣ́мѧ въ зе́млю,
\st{и҆} спи́тъ, и҆ востае́тъ но́щїю и҆ дні́ю: и҆ сѣ́мѧ прозѧба́етъ и҆ расте́тъ, ꙗ҆́коже не вѣ́сть ѻ҆́нъ:
\st{ѿ} себє́ бо землѧ̀ плоди́тъ пре́жде травꙋ̀, пото́мъ кла́съ, та́же и҆сполнѧ́етъ пшени́цꙋ въ кла́сѣ:
\st{є҆}гда́ же созрѣ́етъ пло́дъ, а҆́бїе по́слетъ се́рпъ, ꙗ҆́кѡ наста̀ жа́тва.
\st[]{И҆} гл҃аше: чесомꙋ̀ ᲂу҆подо́бимъ црⷭ҇твїе бж҃їе; и҆лѝ ко́ей при́тчи приложи́мъ є҆̀;
\st{ꙗ҆́}кѡ зе́рно горꙋ́шично, є҆́же є҆гда̀ всѣ́ѧно бꙋ́детъ въ землѝ, мнѣ́е всѣ́хъ сѣ́менъ є҆́сть земны́хъ:
\st{и҆} є҆гда̀ всѣ́ѧно бꙋ́детъ, возраста́етъ, и҆ быва́етъ бо́лѣе всѣ́хъ ѕе́лїй, и҆ твори́тъ вѣ̑тви вє́лїѧ, ꙗ҆́кѡ мощѝ под̾ сѣ́нїю є҆гѡ̀ пти́цамъ небє́снымъ вита́ти.
\st{И҆} таковы́ми при́тчами мно́гими гл҃аше и҆̀мъ сло́во, ꙗ҆́коже можа́хꙋ слы́шати.
\st{Б}ез̾ при́тчи же не гл҃аше и҆̀мъ словесѐ: ѡ҆со́бь же ᲂу҆чн҃кѡ́мъ свои̑мъ сказа́ше всѧ̑.
%
\section % 18
\st[]{И҆} гл҃а и҆̀мъ въ то́й де́нь, ве́черꙋ бы́вшꙋ: пре́йдемъ на ѡ҆́нъ по́лъ.
\st{И҆} ѿпꙋ́щше наро́ды, поѧ́ша є҆го̀ ꙗ҆́коже бѣ̀ въ кораблѝ: и҆ и҆ні́и же корабли̑ бѧ́хꙋ съ ни́мъ.
\st{И҆} бы́сть бꙋ́рѧ вѣ́трена вели́ка: вѡ́лны же влива́хꙋсѧ въ кора́бль, ꙗ҆́кѡ ᲂу҆жѐ погрꙋжа́тисѧ є҆мꙋ̀.
\st{И҆} бѣ̀ са́мъ на кормѣ̀ на возгла́вницѣ спѧ̀. И҆ возбꙋди́ша є҆го̀ и҆ глаго́лаша є҆мꙋ̀: ᲂу҆чт҃лю, не ради́ши ли, ꙗ҆́кѡ погиба́емъ;
\st{И҆} воста́въ запретѝ вѣ́трꙋ и҆ речѐ мо́рю: молчѝ, преста́ни. И҆ ᲂу҆ле́же вѣ́тръ, и҆ бы́сть тишина̀ ве́лїѧ.
\linebreak
\st[]{И҆} речѐ и҆̀мъ: что̀ та́кѡ страшли́ви є҆стѐ; ка́кѡ не и҆́мате вѣ́ры;
\st{И҆} ᲂу҆боѧ́шасѧ стра́хомъ ве́лїимъ и҆ глаго́лахꙋ дрꙋ́гъ ко дрꙋ́гꙋ: кто̀ ᲂу҆̀бо се́й є҆́сть, ꙗ҆́кѡ и҆ вѣ́тръ и҆ мо́ре послꙋ́шаютъ є҆гѡ̀;

% Chapter 5

\vskip-4pt
\section % 19
\hskip0pt%
\St{И҆} прїидо́ша на ѡ҆́нъ по́лъ мо́рѧ, во странꙋ̀ гадари́нскꙋю.
\st{И҆} и҆злѣ́зшꙋ є҆мꙋ̀ и҆з̾ кораблѧ̀, а҆́бїе срѣ́те є҆го̀ ѿ гробѡ́въ человѣ́къ въ дꙋ́сѣ нечи́стѣ,
\st{и҆́}же жили́ще и҆мѧ́ше во гробѣ́хъ, и҆ ни вери́гами никто́же можа́ше є҆го̀ свѧза́ти:
\st{з}анѐ є҆мꙋ̀ мно́гажды пꙋ̑ты и҆ ᲂу҆́жы желѣ̑зны свѧ́занꙋ сꙋ́щꙋ, и҆ растерза́тисѧ ѿ негѡ̀ ᲂу҆́жємъ желѣ̑знымъ и҆ пꙋ́тѡмъ сокрꙋша́тисѧ: и҆ никто́же можа́ше є҆го̀ ᲂу҆мꙋ́чити:
\st{и҆} вы́нꙋ но́щь и҆ де́нь во гробѣ́хъ и҆ въ гора́хъ бѣ̀, вопїѧ̀ и҆ толкі́йсѧ ка́менїемъ.
\st[]{О}у҆зрѣ́въ же і҆и҃са и҆здале́ча, течѐ и҆ поклони́сѧ є҆мꙋ̀,
\st{и҆} возопи́въ гла́сомъ ве́лїимъ, речѐ: что̀ мнѣ̀ и҆ тебѣ̀, і҆и҃се, сн҃е бг҃а вы́шнѧгѡ; заклина́ю тѧ̀ бг҃омъ, не мꙋ́чи менє̀.
\st{Г}л҃аше бо є҆мꙋ̀: и҆зы́ди, дꙋ́ше нечи́стый, ѿ человѣ́ка.
\st{И҆} вопроша́ше є҆го̀: что́ ти є҆́сть и҆́мѧ; И҆ ѿвѣща̀ глаго́лѧ: легеѡ́нъ и҆́мѧ мнѣ̀, ꙗ҆́кѡ мно́зи є҆смы̀.
\st{И҆} моли́ша є҆го̀ мно́гѡ, да не по́слетъ и҆̀хъ внѣ̀ страны̀.
\st{Б}ѣ́ же тꙋ̀ при горѣ̀ ста́до свино́е ве́лїе пасо́мо.
\st[]{И҆} моли́ша є҆го̀ всѝ бѣ́си, глаго́люще: посли́ ны во свинїѧ̑, да въ нѧ̀ вни́демъ.
\st{И҆} повелѣ̀ и҆̀мъ а҆́бїе і҆и҃съ. И҆ и҆зше́дше дꙋ́си нечи́стїи, внидо́ша во свинїѧ̑: и҆ ᲂу҆стреми́сѧ ста́до по бре́гꙋ въ мо́ре, бѧ́хꙋ же ꙗ҆́кѡ двѣ̀ ты́сѧщы, и҆ ᲂу҆топа́хꙋ въ мо́ри.
\st{П}асꙋ́щїи же свинїѧ̑ бѣжа́ша и҆ возвѣсти́ша во гра́дѣ и҆ въ се́лѣхъ. И҆ и҆зыдо́ша ви́дѣти, что̀ є҆́сть бы́вшее.
\st[]{И҆} прїидо́ша ко і҆и҃сови, и҆ ви́дѣша бѣснова́вшагосѧ сѣдѧ́ща и҆ ѡ҆болче́на и҆ смы́слѧща, и҆мѣ́вшаго легеѡ́нъ: и҆ ᲂу҆боѧ́шасѧ.
\st{П}овѣ́даша же и҆̀мъ ви́дѣвшїи, ка́кѡ бы́сть бѣсно́мꙋ, и҆ ѡ҆ свинїѧ́хъ.
\st{И҆} нача́ша моли́ти є҆го̀ ѿитѝ ѿ предѣ̑лъ и҆́хъ.
\st{И҆} влѣ́зшꙋ є҆мꙋ̀ въ кора́бль, молѧ́ше є҆го̀ бѣснова́выйсѧ, дабы̀ бы́лъ съ ни́мъ.
\st{І҆}и҃съ же не дадѐ є҆мꙋ̀, но речѐ є҆мꙋ̀: и҆дѝ въ до́мъ тво́й ко твои̑мъ и҆ возвѣстѝ и҆̀мъ, є҆ли̑ка тѝ гдⷭ҇ь сотворѝ и҆ поми́лова тѧ̀.
\st{И҆} и҆́де и҆ нача̀ проповѣ́дати въ десѧтѝ градѣ́хъ, є҆ли̑ка сотворѝ є҆мꙋ̀ і҆и҃съ. И҆ всѝ дивлѧ́хꙋсѧ.
%
\section % 20
\st[]{И҆} преше́дшꙋ і҆и҃сꙋ въ кораблѝ па́ки на ѡ҆́нъ по́лъ, собра́сѧ наро́дъ мно́гъ ѡ҆ не́мъ. И҆ бѣ̀ при мо́ри.
\st{И҆} сѐ, прїи́де є҆ди́нъ ѿ а҆рхїсѷнагѡ̑гъ, и҆́менемъ і҆аі́ръ, и҆ ви́дѣвъ є҆го̀, падѐ при ногꙋ̀ є҆гѡ̀
\st{и҆} молѧ́ше є҆го̀ мно́гѡ, глаго́лѧ, ꙗ҆́кѡ дщѝ моѧ̀ на кончи́нѣ є҆́сть: да прише́дъ возложи́ши на ню̀ рꙋ́цѣ, ꙗ҆́кѡ да спасе́тсѧ и҆ жива̀ бꙋ́детъ.
\st[]{И҆} и҆́де съ ни́мъ:
%
\section % 21
и҆ по не́мъ и҆дѧ́хꙋ наро́ди мно́зи и҆ ᲂу҆гнета́хꙋ є҆го̀.
\st{И҆} жена̀ нѣ́каѧ сꙋ́щи въ точе́нїи кро́ве лѣ́тъ двана́десѧте,
\st{и҆} мно́гѡ пострада́вши ѿ мнѡ́гъ врачє́въ, и҆ и҆зда́вши своѧ̑ всѧ̑, и҆ ни є҆ди́ныѧ по́льзы ѡ҆брѣ́тши, но па́че въ го́ршее прише́дши:
\st{с}лы́шавши ѡ҆ і҆и҃сѣ, прише́дши въ наро́дѣ созадѝ, прикоснꙋ́сѧ ри́зѣ є҆гѡ̀:
\st{г}лаго́лаше бо, ꙗ҆́кѡ, а҆́ще прикоснꙋ́сѧ ри́замъ є҆гѡ̀, спасе́на бꙋ́дꙋ.
\st{И҆} а҆́бїе и҆зсѧ́кнꙋ и҆сто́чникъ кро́ве є҆ѧ̀: и҆ разꙋмѣ̀ тѣ́ломъ, ꙗ҆́кѡ и҆сцѣлѣ̀ ѿ ра́ны.
\st[]{И҆} а҆́бїе і҆и҃съ разꙋмѣ̀ въ себѣ̀ си́лꙋ и҆зше́дшꙋю ѿ негѡ̀, и҆ ѡ҆бра́щьсѧ въ наро́дѣ, гл҃аше: кто̀ прикоснꙋ́сѧ ри́замъ мои̑мъ;
\st{И҆} глаго́лахꙋ є҆мꙋ̀ ᲂу҆чн҃цы̀ є҆гѡ̀: ви́диши наро́дъ ᲂу҆гнета́ющь тѧ̀, и҆ гл҃еши: кто̀ прикоснꙋ́сѧ мнѣ̀;
\st{И҆} ѡ҆бглѧ́даше ви́дѣти сотво́ршꙋю сїѐ.
\st{Ж}ена́ же ᲂу҆боѧ́вшисѧ и҆ трепе́щꙋщи, вѣ́дѧщи, є҆́же бы́сть є҆́й, прїи́де и҆ припадѐ къ немꙋ̀ и҆ речѐ є҆мꙋ̀ всю̀ и҆́стинꙋ.
\st[]{Ѻ҆́}нъ же речѐ є҆́й: дщѝ, вѣ́ра твоѧ̀ спасе́ тѧ: и҆дѝ въ ми́рѣ и҆ бꙋ́ди цѣла̀ ѿ ра́ны твоеѧ̀.
\st{Є҆}щѐ є҆мꙋ̀ гл҃ющꙋ, и҆ прїидо́ша ѿ а҆рхїсѷнагѡ́га, глаго́люще, ꙗ҆́кѡ дщѝ твоѧ̀ ᲂу҆́мре: что̀ є҆щѐ дви́жеши ᲂу҆чт҃лѧ;
\st{І҆}и҃съ же а҆́бїе слы́шавъ сло́во глаго́лемое, гл҃а а҆рхїсѷнагѡ́гови: не бо́йсѧ, то́кмѡ вѣ́рꙋй.
\st{И҆} не ѡ҆ста́ви по себѣ̀ ни є҆ди́наго и҆тѝ, то́кмѡ петра̀ и҆ і҆а́кѡва и҆ і҆ѡа́нна бра́та і҆а́кѡвлѧ.
\st[]{И҆} прїи́де въ до́мъ а҆рхїсѷнагѡ́говъ и҆ ви́дѣ молвꙋ̀, пла́чꙋщыѧсѧ и҆ крича́щыѧ мно́гѡ.
\st{И҆} вше́дъ гл҃а и҆̀мъ: что̀ мо́лвите и҆ пла́четесѧ; ѻ҆трокови́ца нѣ́сть ᲂу҆мерла̀, но спи́тъ.
\st{И҆} рꙋга́хꙋсѧ є҆мꙋ̀. Ѻ҆́нъ же и҆згна́въ всѧ̑, поѧ́тъ ѻ҆тца̀ ѻ҆трокови́цы и҆ ма́терь, и҆ и҆̀же бѣ́хꙋ съ ни́мъ, и҆ вни́де, и҆дѣ́же бѣ̀ ѻ҆трокови́ца лежа́щи.
\st{И҆} є҆́мь за рꙋ́кꙋ ѻ҆трокови́цꙋ, гл҃а є҆́й: талїѳа̀ кꙋ́мї: є҆́же є҆́сть сказа́емо: дѣви́це, тебѣ̀ гл҃ю, воста́ни.
\st[]{И҆} а҆́бїе воста̀ дѣви́ца и҆ хожда́ше: бѣ́ бо лѣ́тъ двоюна́десѧте. И҆ ᲂу҆жасо́шасѧ ᲂу҆́жасомъ ве́лїимъ.
\st{И҆} запретѝ и҆̀мъ мно́гѡ, да никто́же ᲂу҆вѣ́сть сегѡ̀, и҆ речѐ: дади́те є҆́й ꙗ҆́сти.

\% TODO
